\chapter{VLA 논문 읽기 실전}
\label{ch:paper_reading_practice}

14장에서는 좋은 연구 주제를 고르는 기준을 세웠다. 마지막 장에서는 그 기준을 논문 읽기에 적용한다. 목표는 모든 논문을 같은 깊이로 요약하는 것이 아니라, VLA 연구자가 어떤 순서로 읽고 어떤 산출물을 남겨야 다음 연구 질문으로 이어지는지 정리하는 것이다. 전체 reading catalog는 \cref{app:reading_catalog}에 두고, 본문에는 수업에서 바로 쓸 수 있는 읽기 경로와 worked example을 남긴다.

\section{읽기 순서}

VLA 논문은 최신순으로만 읽으면 구조가 잘 보이지 않는다. 먼저 foundation과 action representation을 읽고, 그 다음 open VLA backbone과 fine-tuning, efficiency, reasoning, benchmark를 읽는 것이 낫다.

\begin{equation}
  \mathrm{Reading\ order}
  =
  \mathrm{Foundation}
  \rightarrow
  \mathrm{Action}
  \rightarrow
  \mathrm{VLA}
  \rightarrow
  \mathrm{System}
  \rightarrow
  \mathrm{Evaluation}.
\end{equation}

이 순서를 따르면 각 논문의 contribution을 ``새 모델''이 아니라 ``새 action interface'', ``새 adaptation recipe'', ``새 system constraint'', ``새 evaluation protocol''로 분해할 수 있다.

\begin{table}[t]
  \centering
  \caption{Compact reading path for a 2024--2026 VLA mini-course}
  \label{tab:compact_reading_path}
  \begin{tabularx}{\textwidth}{>{\bfseries}p{30mm}p{58mm}X}
    \toprule
    Stage & Papers & Expected output \\
    \midrule
    Foundation & SayCan, Diffusion Policy, ACT/ALOHA & Pre-VLA와 VLA의 차이를 controller interface 관점에서 설명한다. \\
    Open VLA & RT-2, Octo, OpenVLA, \(\pi_0\) & 입력, 출력, action space, 학습 데이터, 공개 자원을 표로 비교한다. \\
    Adaptation/action & OpenVLA-OFT, FAST, VQ-VLA, SmolVLA & success, latency, action representation trade-off를 분해한다. \\
    Emerging front & GR00T N1, Gemini Robotics, WorldVLA/VLA-OS, RoboTwin 2.0 & evidence level과 확인 불가능한 claim을 분리한다. \\
    \bottomrule
  \end{tabularx}
\end{table}

\section{4-week reading plan}

대학원 수업이나 랩 세미나에서는 다음 4주 계획을 추천한다. 각 주의 산출물은 짧지만, 모든 논문에 대해 같은 형식을 요구해야 비교가 가능하다.

\begin{table}[t]
  \centering
  \caption{Four-week VLA reading plan}
  \label{tab:four_week_reading_plan}
  \begin{tabularx}{\textwidth}{>{\bfseries}p{14mm}p{56mm}X}
    \toprule
    Week & Papers & Output \\
    \midrule
    1 & SayCan, Diffusion Policy, ACT/ALOHA & Pre-VLA predecessor와 VLA의 경계를 1쪽으로 정리한다. \\
    2 & RT-2, Octo, OpenVLA, \(\pi_0\) & 각 모델의 input-output interface와 action representation 표를 만든다. \\
    3 & OpenVLA-OFT, FAST, VQ-VLA, SmolVLA & action/latency/success trade-off를 claim-source-evidence-caution 표로 쓴다. \\
    4 & GR00T N1, Gemini Robotics, WorldVLA, VLA-OS, RoboTwin 2.0 & evidence ladder를 적용해 peer-reviewed evidence와 company/preprint claim을 분리한다. \\
    \bottomrule
  \end{tabularx}
\end{table}

\section{Worked example 1: OpenVLA-OFT 13-question 적용}

OpenVLA-OFT는 같은 base model에서 fine-tuning recipe가 success와 throughput을 어떻게 바꾸는지 보여 주는 좋은 실전 예제이다 \citep{openvlaoft}. 다음 표는 논문을 수업에서 빠르게 해부하기 위한 압축형 13-question 적용이다.

\begin{table}[t]
  \centering
  \small
  \caption{OpenVLA-OFT 13-question worked example}
  \label{tab:openvla_oft_worked_example}
  \begin{tabularx}{\textwidth}{>{\bfseries}p{8mm}p{24mm}X}
    \toprule
    No. & 항목 & 요약 \\
    \midrule
    01 & 배경 & OpenVLA-style open VLA를 실제 manipulation benchmark와 lab robot에 맞추려면 pretrained backbone만으로는 부족하다. \\
    02 & 문제 & 기존 OpenVLA fine-tuning은 성공률과 action generation throughput 모두에서 deployment 요구를 만족하기 어렵다. \\
    03 & 기존 한계 & Autoregressive action token decoding과 discrete action 표현은 느리고, continuous control에 필요한 정밀도를 잃을 수 있다. \\
    04 & 목표 & 같은 base model을 유지하면서 LIBERO success와 action generation speed를 동시에 개선하는 fine-tuning recipe를 찾는다. \\
    05 & 방법 & Parallel decoding, action chunking, continuous action representation, L1 objective를 결합한 OpenVLA-OFT recipe를 구성한다. \\
    06 & 핵심 아이디어 & VLA 성능은 backbone보다 action interface와 fine-tuning objective가 지배할 수 있으므로, output side를 재설계한다. \\
    07 & 검증 & LIBERO suites와 real-world robot setup에서 기존 OpenVLA-style baseline과 recipe ablation을 비교한다. \\
    08 & 결과 & 논문은 LIBERO 평균 success를 76.5\%에서 97.1\%로 올리고 action generation throughput을 26\(\times\) 개선했다고 보고한다. \\
    09 & 비교 & 비교의 핵심은 같은 base model 통제 아래 action representation, decoding, objective를 바꾼 ablation을 보는 것이다. \\
    10 & 의의 & Open VLA를 실제 lab task에 맞추는 병목이 단순 model size가 아니라 action decoding과 adaptation recipe임을 보여 준다. \\
    11 & 한계 & 결과는 benchmark protocol, hardware, controller, LIBERO split 조건에 묶이므로 다른 robot deployment로 바로 일반화하면 안 된다. \\
    12 & 향후 과제 & Real-time safety layer, broader robot platforms, failure recovery data, cross-benchmark transfer로 recipe의 범위를 검증해야 한다. \\
    13 & 자원 공개 & OpenVLA의 공개 backbone과 함께 읽되, OFT의 code/checkpoint 공개 범위와 재현 가능한 evaluation script를 별도로 확인해야 한다. \\
    \bottomrule
  \end{tabularx}
\end{table}

\section{Worked example 2: WorldVLA/VLA-OS 13-question 적용}

Reasoning과 world model 논문은 자칫 그럴듯한 개념 설명으로 끝나기 쉽다. WorldVLA와 VLA-OS는 ``미래 예측이나 planning representation이 실제 action-level utility를 높이는가''라는 질문으로 읽어야 한다 \citep{worldvla,vlaos}.

\begin{table}[t]
  \centering
  \small
  \caption{WorldVLA/VLA-OS 13-question worked example}
  \label{tab:worldvla_vlaos_worked_example}
  \begin{tabularx}{\textwidth}{>{\bfseries}p{8mm}p{24mm}X}
    \toprule
    No. & 항목 & 요약 \\
    \midrule
    01 & 배경 & Long-horizon VLA는 현재 관측에서 바로 action을 내는 reactive policy만으로는 planning, recovery, safety를 충분히 다루기 어렵다. \\
    02 & 문제 & World model 또는 planning representation이 실제 행동 선택에 도움을 주는지, 단순 text reasoning이나 pixel prediction을 넘어 검증해야 한다. \\
    03 & 기존 한계 & 많은 reasoning-style VLA는 중간 문장이나 plan의 자연스러움은 보이지만 closed-loop utility와 failure recovery evidence가 약하다. \\
    04 & 목표 & Action-conditioned future prediction 또는 structured planning representation이 candidate ranking, recovery, task success를 개선하는지 확인한다. \\
    05 & 방법 & WorldVLA는 action과 visual generation/understanding을 묶는 autoregressive world model 방향이고, VLA-OS는 planning representation과 paradigm을 구조적으로 비교한다. \\
    06 & 핵심 아이디어 & Reasoning의 가치는 설명문이 아니라 action 후보를 더 잘 고르고 실패를 더 빨리 감지하게 만드는 중간 변수에서 나온다. \\
    07 & 검증 & Prediction metric, candidate ranking, planning ablation, closed-loop 또는 benchmark task success를 분리해 확인해야 한다. \\
    08 & 결과 & 핵심 결과는 pixel score가 아니라 world model 또는 planning representation을 넣었을 때 success, recovery, safety metric이 좋아지는지이다. \\
    09 & 비교 & No-world-model, no-planning, direct policy, text-only reasoning baseline과의 통제 비교가 필요하다. \\
    10 & 의의 & 검증이 충분하면 VLA를 reactive policy에서 prediction-aware policy로 확장하는 근거가 된다. \\
    11 & 한계 & 미래 예측은 contact dynamics, occlusion, long horizon에서 빠르게 틀어질 수 있고, rollout latency가 control loop를 망칠 수 있다. \\
    12 & 향후 과제 & Short-horizon verification, uncertainty-aware replanning, real-robot recovery benchmark, safety-critical counterfactual test가 필요하다. \\
    13 & 자원 공개 & Preprint/project report의 공개 코드, benchmark script, model checkpoint, evaluation protocol을 확인하고 없으면 claim-level evidence로 낮춰 읽는다. \\
    \bottomrule
  \end{tabularx}
\end{table}

\section{Worked example 3: Humanoid and closed industrial VLA evidence split}

GR00T N1, Helix, Gemini Robotics를 한 표에 놓으면 VLA의 최신 흐름은 잘 보이지만, evidence level을 구분하지 않으면 hype와 검증이 섞인다. 이 worked example의 목적은 open arXiv source, company report, closed VLA, embodied reasoning model을 같은 기준으로 분리하는 것이다 \citep{grootn1,helix,geminirobotics,geminiroboticser}.

\begin{table}[t]
  \centering
  \small
  \caption{Evidence-ladder worked example for humanoid and industrial VLA}
  \label{tab:industrial_vla_evidence_worked_example}
  \begin{tabularx}{\textwidth}{>{\bfseries}p{26mm}p{18mm}p{34mm}X}
    \toprule
    Source & Level & What the source claims & How to read it \\
    \midrule
    GR00T N1 & B & Humanoid foundation model with vision-language module and diffusion transformer motor-action module. & Open foundation claim은 중요하지만, 공개 artifact, data mixture, simulation/real split, humanoid platform 조건을 따로 확인한다. \\
    Helix & C & Generalist humanoid VLA for full-upper-body control in an industrial product setting. & Company report/demo evidence로 분류하고, benchmark score나 reproducible paper evidence처럼 비교하지 않는다. \\
    Gemini Robotics 1.5 & C & Closed industrial VLA that maps visual information and instruction to robot motor commands. & Closed model이므로 model weight, controller, data, safety layer의 비공개 범위를 먼저 표시한다. \\
    Gemini Robotics-ER 1.6 & C & Embodied reasoning model for robotics rather than the same object as the action VLA. & Reasoning model과 motor policy를 구분하고, reasoning claim이 recovery/safety/action metric으로 닫히는지 본다. \\
    Reviewer question & -- & 이 source가 field direction을 보여 주는가, 아니면 reproducible technical evidence를 제공하는가? & Trend discussion에는 쓸 수 있지만, SOTA claim에는 evidence level을 명시해야 한다. \\
    \bottomrule
  \end{tabularx}
\end{table}

\section{논문 읽기 노트 템플릿}

매일 여러 편을 읽을 때는 13-question 전체를 항상 쓰기보다 짧은 note template으로 시작하는 것이 효율적이다.

\begin{itemize}
  \item 이 논문의 정확한 병목:
  \item 제안한 mechanism:
  \item 가장 믿을 만한 evidence:
  \item 가장 약한 baseline 또는 confound:
  \item 내 연구로 가져올 수 있는 지점:
  \item 실제 deployment에서 남는 위험:
\end{itemize}

이 template은 13-question summary보다 짧다. 중요한 논문만 13-question으로 확장하고, 나머지는 이 노트로 cataloging하면 field map을 빠르게 만들 수 있다.

\section{Assignment difficulty tags}

각 장의 reading assignment는 다음 난이도 태그 중 하나로 읽는다. 수업에서는 Basic을 개인 과제로, Intermediate를 세미나 토론 과제로, Advanced와 Capstone을 프로젝트 제안서로 연결하면 좋다.

\begin{table}[t]
  \centering
  \caption{Assignment difficulty tags used in the book}
  \label{tab:assignment_difficulty_tags}
  \begin{tabularx}{\textwidth}{>{\bfseries}p{28mm}X}
    \toprule
    Tag & Expected work \\
    \midrule
    Basic & 13-question summary, input-output interface, key claim extraction처럼 논문 독해의 기본 구조를 채운다. \\
    Intermediate & Action space, benchmark tuple, claim-source-evidence-caution table처럼 비교 가능한 분석 표를 만든다. \\
    Advanced & Ablation, failure taxonomy, safety checklist, deployment readiness처럼 논문 claim을 재설계하거나 검증 protocol을 제안한다. \\
    Capstone & 여러 source의 evidence level을 비교해 새로운 VLA 연구 병목과 실험 계획을 제안한다. \\
    \bottomrule
  \end{tabularx}
\end{table}

\section{VLA 연구자가 가져가야 할 7개 원칙}

이 책의 마지막 결론은 다음 일곱 문장으로 요약할 수 있다.

\begin{enumerate}
  \item 모델명보다 interface를 먼저 보라.
  \item Action representation은 구현 세부사항이 아니라 핵심 contribution이다.
  \item Success rate는 benchmark tuple에 조건부인 숫자다.
  \item Real-time constraint 없이는 robot policy claim이 닫히지 않는다.
  \item Reasoning은 긴 text가 아니라 recovery, uncertainty, safety로 검증하라.
  \item Closed industrial claim은 evidence ladder로 분리하라.
  \item 좋은 VLA 연구 주제는 bottleneck, mechanism, evidence가 한 문장으로 연결될 때 나온다.
\end{enumerate}

VLA 분야는 빠르게 변하지만, 좋은 논문을 읽는 기준은 크게 변하지 않는다. 입력, 출력, action space, controller interface, learning signal, evaluation protocol, source status를 확인하면 hype와 실제 engineering contribution을 분리할 수 있다. 이 책의 목적도 바로 그 분리 능력을 만드는 데 있다.

\begin{assignmentbox}{Final assignment [Capstone]}
관심 있는 VLA 논문 한 편을 골라 6-line note, 13-question summary, claim-source-evidence-caution table을 모두 작성하고, 마지막에 그 논문이 새로운 연구 주제로 이어지는 병목 하나를 제안하라.
\end{assignmentbox}
